%%%%%%%%%%%%%%%%%%%%%%%%%%%%%%%%%%%%%%%%%
% FRI Data Science_report LaTeX Template
% Version 1.0 (28/1/2020)
% 
% Jure Demšar (jure.demsar@fri.uni-lj.si)
%
% Based on MicromouseSymp article template by:
% Mathias Legrand (legrand.mathias@gmail.com) 
% With extensive modifications by:
% Antonio Valente (antonio.luis.valente@gmail.com)
%
% License:
% CC BY-NC-SA 3.0 (http://creativecommons.org/licenses/by-nc-sa/3.0/)
%
%%%%%%%%%%%%%%%%%%%%%%%%%%%%%%%%%%%%%%%%%


%----------------------------------------------------------------------------------------
%	PACKAGES AND OTHER DOCUMENT CONFIGURATIONS
%----------------------------------------------------------------------------------------
\documentclass[fleqn,moreauthors,10pt]{ds_report}
\usepackage[english]{babel}

\graphicspath{{fig/}}




%----------------------------------------------------------------------------------------
%	ARTICLE INFORMATION
%----------------------------------------------------------------------------------------

% Header
\JournalInfo{FRI Natural language processing course 2021}

% Interim or final report
\Archive{Project report} 
%\Archive{Final report} 

% Article title
\PaperTitle{Natural language processing project proposal} 

% Authors (student competitors) and their info
\Authors{Strahinja Đorđević, Stefan Anđelković}

% Advisors
\affiliation{\textit{Advisors: Slavko Žitnik}}

% Keywords
\Keywords{Keyword1, Keyword2, Keyword3 ...}
\newcommand{\keywordname}{Keywords}


%----------------------------------------------------------------------------------------
%	ABSTRACT
%----------------------------------------------------------------------------------------

\Abstract{
The abstract goes here. The abstract goes here. The abstract goes here. The abstract goes here. The abstract goes here. The abstract goes here. The abstract goes here. The abstract goes here. The abstract goes here. The abstract goes here. The abstract goes here. The abstract goes here. The abstract goes here. The abstract goes here. The abstract goes here. The abstract goes here. The abstract goes here. The abstract goes here. The abstract goes here. The abstract goes here. The abstract goes here. The abstract goes here. The abstract goes here. The abstract goes here. The abstract goes here. The abstract goes here.
}

%----------------------------------------------------------------------------------------

\begin{document}

% Makes all text pages the same height
\flushbottom 

% Print the title and abstract box
\maketitle 

% Removes page numbering from the first page
\thispagestyle{empty} 

%----------------------------------------------------------------------------------------
%	ARTICLE CONTENTS
%----------------------------------------------------------------------------------------

\section*{Introduction}
    Since the popularization of large language models, there have been many solutions developed to solve different every day problems. One such example are chat bots, the most popular being ChatGPT. Many solutions take a pretrained model and finetune them for specialized use-cases, usually in the form of an assistant that the user can ask question to (e.g. an online shopping website, where you can ask questions about products). This paper will present one such use-case.\\

    At the Faculty of computer and information science, University of Ljubljana (slo. FRI), a common problem students face is picking elective courses. This seems to be an especially relevant problem at FRI because it has so many different courses. For this reason, we propose a chat assistant that students can ask questions about courses, to help them decide which courses to pick. Students can say what topics interest them, how many ECTS credits they need, etc., and the chatbot will give them suggestions for courses that best meet their needs. We also plan on implementing a mechanism that will let the student give specific time slot requirements. This way students could give the bot their time table (maybe in the form of their student ID), and request that the bot avoids any time slot overlap with their current courses.

\subsection*{Related work}

Specialized AI assistants have emerged across various industries. In finance, Bank of America's Erica helps nearly 50 million users with spending management and financial guidance. In retail, Bindly offers personalized shopping recommendations through natural conversation, while H\&M's chatbot provides outfit and style advice on messaging platforms. In the culinary domain, the Cookidoo Assistant integrated into Bimby TM7 helps users discover recipes based on available ingredients and can generate new preparations tailored to dietary constraints.

In recent years, the use of chatbots has grown rapidly across various fields, including education. Ranoliya et al.~\cite{ranoliya2017chatbot} developed a chatbot for university-related FAQs, while Clarizia et al.~\cite{clarizia2018chatbot} created an educational chatbot prototype employing natural language processing and domain ontologies to effectively answer student queries. Most notably, a recent comprehensive survey by Swacha and Gracel~\cite{swacha2025rag} analyzed 47 papers on Retrieval-Augmented Generation (RAG) chatbots in education, examining their purposes, thematic scope, and evaluation methods.

At the Faculty of Computer and Information Science, students face difficulty selecting elective courses due to the large number of options. We propose a chat assistant that helps students choose courses based on their interests, ECTS requirements, and scheduling constraints. Such a chatbot would greatly benefit the FRI faculty by reducing the administrative burden on staff, who currently spend significant time answering individual student inquiries about course selection, while simultaneously helping students make more informed decisions about their academic paths.

\subsection*{Provisional dataset}
The main source of information our chatbot will use are PDF files containing in-detph information about every course that the Faculty offers \cite{uni_fri_zbornik, vss_fri_zbornik}. These contain general information about each course, the number of credits they give, the semester and year they are held in, the field of study the course is meant for, as well as some other details. We will also try to obtain information about the time table for the current year \cite{urnik_fri}. We will also include all generally available data on the FRI website \cite{ul_fri}.

%------------------------------------------------

\section*{Methods}

\subsection*{Document Processing and Indexing}

The knowledge base for our chatbot is built from PDF course catalogues for the University of Ljubljana, Faculty of Computer and Information Science. These documents are first converted to Markdown format using the \texttt{pymupdf4llm} library, which preserves tabular structures and headings relevant to course descriptions. The resulting Markdown files are stored locally and serve as the primary input for the retrieval pipeline.

For indexing, documents are split into overlapping chunks using LangChain's \\ \texttt{RecursiveCharacterTextSplitter} with a chunk size of 1000 characters and an overlap of 150 characters. Each chunk is then embedded using the \\ \texttt{sentence-transformers/all-MiniLM-L6-v2} \\ model, which produces L2-normalised 384-dimensional dense vectors. The resulting embeddings are stored in a FAISS vector index that is persisted to disk and can be reloaded between sessions.

\subsection*{Retrieval}

At query time, the user's question is first preprocessed by expanding domain-specific abbreviations (e.g., \textit{uni} $\rightarrow$ \textit{univerzitetni program}, \textit{ects} $\rightarrow$ \textit{kreditne točke}). The processed query is then embedded with the same model used during indexing and used to retrieve the $k$ most relevant chunks from the FAISS index. We support two retrieval strategies: standard cosine similarity search and Maximal Marginal Relevance (MMR), which balances relevance and diversity among retrieved chunks.

The value of $k$ is determined dynamically. Before each query, a dedicated model call extracts any explicit quantity mentioned in the question (e.g., \textit{``give me 3 courses''}); if a number is found, $k$ is scaled accordingly. This allows the system to adapt the size of the retrieved context to the user's informational need.

\subsection*{Generation}

We use Meta's \texttt{Llama-3-8B-Instruct} model as the generative backbone, loaded in 4-bit NormalFloat quantisation (NF4) with double quantisation enabled via BitsAndBytes, allowing inference on a single consumer GPU. The model is wrapped in a LangChain \texttt{HuggingFacePipeline} with greedy decoding (\texttt{do\_sample=False}) and a repetition penalty of 1.1 to reduce degenerate outputs.

Retrieved chunks are concatenated and injected into a structured system prompt that instructs the model to answer exclusively from the provided context and to always include the course name, ECTS credits, and semester when listing courses. The prompt is formatted using the model's built-in chat template via \texttt{apply\_chat\_template}.

\subsection*{Conversational Memory}

The system supports multi-turn dialogue by maintaining a conversation history as a list of question–answer pairs. This history is serialised to a JSON file after each turn and reloaded at the start of the next session, ensuring continuity across separate invocations of the pipeline. The full history is prepended to each new prompt so the model can resolve references to previous turns.

\subsection*{RAG Pipeline}

The complete pipeline is implemented using LangChain's \texttt{LCEL} (LangChain Expression Language) and consists of four stages: (1) context retrieval, (2) prompt construction, (3) language model generation, and (4) output parsing. The pipeline accepts a dictionary containing the current question and conversation history, and returns the model's answer as a plain string.

%------------------------------------------------

\section*{Results}

To evaluate the performance of the proposed RAG-based chatbot, we conducted a series of test queries covering different levels of complexity. The queries ranged from simple factual lookups to more detailed, multi-course descriptions. Below we present and analyse the results of four representative test cases.

\subsection*{Simple Factual Queries}

For straightforward queries, the system performed reliably. When asked to name two courses related to embedded systems, the chatbot correctly identified \textit{Vgrajeni sistemi} (Embedded Systems) and a second course involving PIC microcontrollers. Similarly, when asked about courses covering machine learning and artificial intelligence, the system returned a list of eight relevant courses, including \textit{Strojno učenje}, \textit{Umetna inteligencija}, \textit{Računalniška vidnost}, and \textit{Spodbujevalno učenje}. While all of these courses exist, the system does make up the names of some courses (eg. "Računalniška vidnost", with the real name of the course being "Umetno zaznavanje).

These results demonstrate that for well-scoped queries with a clear answer present in the retrieved context, the pipeline retrieves appropriate chunks and the language model produces accurate, grounded responses.

\subsection*{Multi-Course Detail Queries}

When the user requested more detail about two specific courses --- \textit{Deep Learning} and \textit{Computer Vision} --- the system produced structured descriptions for both, including topics covered and learning goals. While the high-level structure of the responses was coherent and the course names were correctly identified, we observed the first signs of hallucination at this level of complexity. Specifically, the course \textit{Gladko učenje} was listed as the Slovenian title for Deep Learning, which is incorrect; the accurate translation is \textit{Globoko učenje}. Additionally, some of the listed subtopics (e.g., specific frameworks and techniques) could not be verified against the source documents, suggesting the model partially generated content beyond what the retrieved context contained.

\subsection*{Complex Thematic Queries}

For more open-ended queries, such as asking for all courses related to software architectures and requesting detailed descriptions of each, the quality of responses degraded noticeably. While the system did identify two plausible courses --- \textit{Architecture Design} and \textit{Implementation and Advanced Topics in Software Engineering} --- the detailed content attributed to these courses included topics that do not appear in the source PDFs, such as specific architectural patterns (MVC, EDA, pipe and filter) and methodological frameworks (Kruchten's 4+1), which were either absent from or only briefly mentioned in the retrieved chunks.

This behaviour is consistent with a known failure mode of RAG systems: when the retrieved context is insufficient to fully answer a complex query, the language model falls back on its parametric knowledge and generates plausible-sounding but unverified information. In a practical deployment, this is particularly problematic because the fabricated course details could mislead students during course selection.

\subsection*{Summary}

Overall, the system shows promising performance for simple retrieval tasks, but struggles with complex queries that require synthesising information across many documents or generating detailed descriptions not explicitly present in the source material.
%------------------------------------------------

\section*{Discussion}

Use the Discussion section to objectively evaluate your work, do not just put praise on everything you did, be critical and exposes flaws and weaknesses of your solution. You can also explain what you would do differently if you would be able to start again and what upgrades could be done on the project in the future.


%------------------------------------------------

\section*{Acknowledgments}

Here you can thank other persons (advisors, colleagues ...) that contributed to the successful completion of your project.


%----------------------------------------------------------------------------------------
%	REFERENCE LIST
%----------------------------------------------------------------------------------------
\bibliographystyle{unsrt}
\bibliography{report}


\end{document}
